\section{Results}
\label{sec:results}

% Style note (v2, 2026-05-07): scaffolded in the EMNLP 2025 Outstanding-paper
% Setup/Hypothesis/Results/Discussion micro-template. Each sub-section
% preregisters its hypothesis before reporting findings. Numerical
% [TBD] placeholders mark cells that the Phase B/C/F runs will populate;
% the rhetorical scaffold and statistical protocol are settled here so
% that camera-ready edits are local.

We organize the results around three questions, each tested over the same main matrix with the protocol of \S\ref{sec:experimental-setup}.

\subsection{Q1: Does whole-pool DualEnd preserve effectiveness while approximately halving serial LLM calls?}
\label{sec:results-cost-quality}

\paragraph{Setup.}
We compare the seven methods within each (model, dataset, pool-size) cell at the matched operating point $W{=}k{=}N$. The within-MaxContext contrast is between \textsc{DualEnd} and either single-extreme MaxContext variant, where call counts are deterministic (Table~\ref{tab:call-accounting}). The cross-family contrast is between the three MaxContext variants and the four standard windowed setwise baselines.

\paragraph{Hypothesis.}
We predict that \textsc{DualEnd} is non-inferior to
\textsc{MaxContext-TopDown} and \textsc{MaxContext-BottomUp} on per-query nDCG@10 at margin $\delta{=}0.01$, on at least 80\% of cells (the H3 threshold). We further predict that the three MaxContext variants are competitive with the standard windowed setwise baselines at the same matched operating point, since context-length is no longer the binding constraint that motivated the windowed design. We do \emph{not} predict that \textsc{DualEnd} dominates the field; the central efficiency claim is about serial LLM calls, not effectiveness.

\paragraph{Results.}
Table~\ref{tab:main-deltas} reports per-cell signed nDCG@10 differences $\Delta_{\mathrm{TD}} = \mathrm{nDCG@10}(\textsc{DualEnd}) - \mathrm{nDCG@10}(\textsc{MaxContext-TopDown})$ and $\Delta_{\mathrm{BU}} = \mathrm{nDCG@10}(\textsc{DualEnd}) - \mathrm{nDCG@10}(\textsc{MaxContext-BottomUp})$ at $N{=}50$, with $\checkmark$ / $\blacktriangle$ / $\times$ markers for the H1 non-inferiority verdict. [TBD: per-cell verdicts after Phase B completion.] Aggregated across the 4-model main-paper subset and 6 datasets, [TBD: $X/24$ cells favor non-inferiority for $\Delta_{\mathrm{TD}}$ and $Y/24$ for $\Delta_{\mathrm{BU}}$ under Holm-Bonferroni at $\alpha{=}0.05$]. Figure~\ref{fig:pareto} plots the mean LLM calls vs.\ mean nDCG@10 across the pool sweep $N \in \{10,20,30,40,50,100\}$.

% Table 2 — Compressed signed-Δ matrix at N=50.
% PLACEHOLDER: this table is camera-ready-only. Once Phase B finishes,
% replace this file with the populated 4×4×2 matrix:
%   Rows = 4 Tier-1 models (Qwen3.5-4B, Qwen3.5-9B, Llama-3.1-8B-Instruct, Ministral-3-8B-Instruct-2512)
%   Column blocks = {DL19, DL20, BEIR-mean, BEIR-min}
%   Per-block columns = Δ_TD = nDCG@10(DualEnd) − nDCG@10(TopDown),
%                       Δ_BU = nDCG@10(DualEnd) − nDCG@10(BottomUp)
%   Cell entry = signed Δ with marker ✓ / ▲ / ✗ for the H1 verdict.

\begin{table*}[t]
\centering
\small
\caption{Per-cell signed nDCG@10 differences at $N{=}50$ between \textsc{MaxContext-DualEnd} and the two single-extreme controls, with $\checkmark$ / $\blacktriangle$ / $\times$ markers for the per-cell H1 non-inferiority verdict (one-sided paired permutation, $\delta{=}0.01$, Holm-Bonferroni at $\alpha{=}0.05$). \textit{Placeholder: populated at camera-ready from Phase B.}}
\label{tab:main-deltas}
\begin{tabular}{l cc cc cc cc}
\toprule
& \multicolumn{2}{c}{\textbf{DL19}} & \multicolumn{2}{c}{\textbf{DL20}} & \multicolumn{2}{c}{\textbf{BEIR-mean}} & \multicolumn{2}{c}{\textbf{BEIR-min}} \\
\cmidrule(lr){2-3} \cmidrule(lr){4-5} \cmidrule(lr){6-7} \cmidrule(lr){8-9}
\textbf{Model} & $\Delta_{\mathrm{TD}}$ & $\Delta_{\mathrm{BU}}$ & $\Delta_{\mathrm{TD}}$ & $\Delta_{\mathrm{BU}}$ & $\Delta_{\mathrm{TD}}$ & $\Delta_{\mathrm{BU}}$ & $\Delta_{\mathrm{TD}}$ & $\Delta_{\mathrm{BU}}$ \\
\midrule
Qwen3.5-4B & --- & --- & --- & --- & --- & --- & --- & --- \\
Qwen3.5-9B & --- & --- & --- & --- & --- & --- & --- & --- \\
Llama-3.1-8B-Instruct & --- & --- & --- & --- & --- & --- & --- & --- \\
Ministral-3-8B-Instruct-2512 & --- & --- & --- & --- & --- & --- & --- & --- \\
\bottomrule
\end{tabular}
\end{table*}

% Figure 2 — Pareto plot of mean LLM calls vs. mean nDCG@10.
% PLACEHOLDER: this figure is camera-ready-only. Once Phase B + Phase C
% data lands, replace this file with the rendered Pareto plot.
% Top panel: DL19+DL20 mean across N ∈ {10..100}.
% Bottom panel: BEIR-mean across N ∈ {10, 30, 50, 100}.
% Three line series per panel: DualEnd, TopDown, BottomUp.

\begin{figure}[t]
\centering
\framebox[\columnwidth][c]{\rule{0pt}{4.0cm}\textit{[TBD: Pareto plot, top = DL19+DL20 mean, bottom = BEIR-mean]}}
\caption{Pareto plot of mean LLM calls (x-axis) versus mean nDCG@10 (y-axis) across pool sizes $N \in \{10,20,30,40,50,100\}$. Top panel reports DL19+DL20 mean; bottom panel reports BEIR-mean over the four Tier-1 BEIR domains. Three line series: \textsc{MaxContext-DualEnd}, \textsc{MaxContext-TopDown}, \textsc{MaxContext-BottomUp}. The bottom panel uses only $N \in \{10, 30, 50, 100\}$ to match the experiment-plan trim. \textit{Placeholder: rendered version replaces this stub at camera-ready.}}
\label{fig:pareto}
\end{figure}


\paragraph{Discussion.}
[TBD interpretation paragraph.] If the verdict pattern is supportive, this is direct evidence that joint best--worst elicitation amortizes the serial-call cost of whole-pool setwise selection. If supportive only below $N{=}100$, the discussion frames a long-context saturation point rather than a flat win. If the verdict is mixed, the call-count reduction remains exact (Table~\ref{tab:call-accounting}) but its practical value depends on the model family --- a discussion we resume in Q2.

\subsection{Q2: Does the result generalize across model families?}
\label{sec:results-family}

\paragraph{Setup.}
Each (method, dataset, pool-size) verdict from Q1 is broken out by model family --- five Qwen3.5 checkpoints, one Llama-3.1, three Ministral-3 --- and by parameter scale within Qwen3.5 (0.8B--27B). Tier-1 main-paper cells use four checkpoints chosen to vary family while holding scale near the 8B class; the remaining cells populate the appendix.

\paragraph{Hypothesis.}
We predict that the within-family DualEnd-vs-single-extreme contrast is qualitatively consistent across the three families, even where its magnitude is not. We do not predict that all families behave identically; differences in chat template, tokenizer, and long-context implementation are expected to produce family-specific effect sizes.

\paragraph{Results.}
[TBD per-family signed-$\Delta$ summary across the four Tier-1 checkpoints and six datasets at $N{=}50$, with per-family verdict counts.] Figure~\ref{fig:pareto} bottom panel reports the BEIR-mean Pareto curve per family. Family-stratified scale behavior within Qwen3.5 (0.8B--27B) is reported in the appendix.

\paragraph{Discussion.}
[TBD interpretation: whether the finding is a property of the algorithm or of a specific family.] The cross-family design is the principal guard against the single-family confound that motivated this study~(\S\ref{sec:methodology}); the specific pattern of agreement or disagreement across Qwen3.5, Llama-3.1, and Ministral-3 is the cleanest evidence we can report on family-generality.

\subsection{Q3: How sensitive is whole-pool MaxContext to candidate order?}
\label{sec:results-position-bias}

\paragraph{Setup.}
We pair each MaxContext-only forward run from Q1 with two per-comparison position-bias controls applied to the live pool $L$ before each LLM call: a deterministic \textbf{reverse} of $L$, and a fixed-seed shuffle of $L$ (seed 929). The required matrix covers three representative models, four representative datasets, three MaxContext methods, six pool sizes, and two control conditions ($3 \times 4 \times 3 \times 6 \times 2 = 432$ runs; \S\ref{sec:experimental-setup}). The analysis reports paired within-query nDCG@10 deltas for forward vs.\ reverse and forward vs.\ shuffle, plus selected-label position frequencies.

\paragraph{Hypothesis.}
We predict that whole-pool MaxContext is \emph{not} order-invariant: the forward-vs-reverse and forward-vs-shuffle deltas should be detectable above the deterministic-decoding noise floor, with magnitudes that correlate with pool size $N$ and may correlate with model family. We take this as a property of long-context LLM prompting~\citep{liu2024lostmiddle,tang2024foundmiddle,zheng2024mcqbias,shi2025positionbias,pradeep2023rankzephyr} rather than a critique of MaxContext.

\paragraph{Results.}
[TBD: per-cell paired delta distribution; aggregated summary statistics (mean, SD, range, worst-pair $|\Delta|$); selected-label position frequencies.] We expect the position-frequency view to expose whether selected-label distributions concentrate at specific live-pool positions under forward order and shift predictably under reverse and shuffle.

\paragraph{Discussion.}
[TBD interpretation.] If forward-vs-reverse deltas are large and systematic, this is a position-bias finding for whole-pool MaxContext that windowed setwise reranking does not surface --- the finding is specific to the long-context regime and does not refute the call-count result in Q1. If deltas are small, the long-context regime tolerates pool reordering and the call-count result is robust. The deterministic 10$\times$ stability check (\S\ref{sec:experimental-setup}) is reported separately as a byte-equality guard rather than as input-order evidence.
